\begin{proofbox}
\begin{center}
\begin{minipage}{0.98\linewidth}\centering
\textbf{两块正方形比两块矩形，多出一块差的平方}\par\medskip
\begin{tikzpicture}[font=\small]
\filldraw[fill=CProof!10,draw=CProof,thick] (0,0) rectangle (2,2);
\node at (1,1) {$u^2$};\node[below] at (1,0) {$u$};
\filldraw[fill=CProof!10,draw=CProof,thick] (2.6,0) rectangle (3.6,1);
\node at (3.1,0.5) {$v^2$};\node[below] at (3.1,0) {$v$};
\node at (4.3,1) {$=$};
\filldraw[fill=CExample!12,draw=CExample,thick] (5,0) rectangle (7,1);
\filldraw[fill=CExample!12,draw=CExample,thick] (5,1.4) rectangle (7,2.4);
\node at (6,0.5) {$uv$};\node at (6,1.9) {$uv$};
\node at (7.5,1) {$+$};
\filldraw[fill=CTrap!12,draw=CTrap,thick] (8,0.5) rectangle (9,1.5);
\node[above] at (8.5,1.5) {$(u-v)^2$};
\end{tikzpicture}\par\smallskip
图示 $u=2,v=1$；一般取 $u=\sqrt{Ax},v=\sqrt{B/x}$（$x>0$）。
面积恒等式 $u^2+v^2=2uv+(u-v)^2$ 表明：$Ax+B/x$ 比 $2\sqrt{AB}$ 多出的量总非负；仅 $u=v$ 时为零。
\end{minipage}\end{center}


	\textbf{\textcolor{CProof}{思路提示}}
	对 $x>0$ 部分直接用均值不等式；对 $x<0$ 部分通过换元 $x=-t$ 化为正数情形，再利用不等式方向反转得到上界。

	\medskip
	\textbf{\textcolor{CProof}{正式推导}}
	\begin{description}[style=nextline, leftmargin=2.8em, labelsep=0.5em, itemsep=0.55em]
		\item[\textbf{步骤一（正数最小值）：}]
			证明：当 $x>0$ 时，$Ax>0,\ \frac{B}{x}>0$，由均值不等式
			\[
				\begin{aligned}
					Ax+\frac{B}{x} &\ge 2\sqrt{Ax\cdot\frac{B}{x}} \\ &= 2\sqrt{AB}.
			\end{aligned}
			\]
			\par\textbf{理由：} 均值不等式 $a+b\ge 2\sqrt{ab}$（$a,b>0$）立得。等号成立当且仅当 $Ax=\frac{B}{x}$，即 $x=\sqrt{\frac{B}{A}}$（$x>0$）.
		\item[\textbf{步骤二（负数最大值）：}]
			考虑 $x<0$，令 $x=-t$（$t>0$），则
			\[
				\begin{aligned}
					f(x) &= A(-t)+\frac{B}{-t} \\ &= -\left(At+\frac{B}{t}\right).
			\end{aligned}
			\]
			由步骤一，$t>0$ 时有 $At+\frac{B}{t}\ge 2\sqrt{AB}$，因此
			\[
				\begin{aligned}
					f(x) &= -\left(At+\frac{B}{t}\right) \\ &\le -2\sqrt{AB}.
			\end{aligned}
			\]
			\par\textbf{理由：} 通过 $x=-t$ 将 $x<0$ 转化为 $t>0$，然后套用正数范围的不等式，注意不等号方向因取负而反转。等号成立当 $t=\sqrt{\frac{B}{A}}$，即 $x=-\sqrt{\frac{B}{A}}$.
		\item[\textbf{步骤三（值域综合）：}]
			由步骤一知，$x>0$ 时 $f(x)\ge 2\sqrt{AB}$，且当 $x\to 0^+$ 时 $\frac{B}{x}\to +\infty$，当 $x\to +\infty$ 时 $Ax\to +\infty$，故 $f(x)$ 在 $(0,+\infty)$ 上可取得 $[2\sqrt{AB},+\infty)$ 中的所有值。类似地，由步骤二知，$x<0$ 时 $f(x)\le -2\sqrt{AB}$，且 $x\to 0^-$ 时 $f(x)\to -\infty$，$x\to -\infty$ 时 $f(x)\to -\infty$，故 $f(x)$ 在 $(-\infty,0)$ 上可取得 $(-\infty,-2\sqrt{AB}]$ 中的所有值。因此综合值域为
			\[
				(-\infty,-2\sqrt{AB}\;]\cup[\;2\sqrt{AB},+\infty).
			\]
			\par\textbf{理由：} 由连续性及极限趋势可知 $f(x)$ 在两个区间内都能取到所有中间值（利用介值性），将两段值域取并集即得.
	\end{description}

	\medskip
	\textbf{\textcolor{CProof}{结论回扣}}
	至此，我们严格证明了 $f(x)=Ax+\frac{B}{x}$（$A,B>0$）在 $x>0$ 时最小值为 $2\sqrt{AB}$，在 $x<0$ 时最大值为 $-2\sqrt{AB}$，并得到值域为 $(-\infty,-2\sqrt{AB}]\cup[2\sqrt{AB},+\infty)$. 同时给出了取等条件 $x=\pm\sqrt{\frac{B}{A}}$.


\medskip
\textbf{\textcolor{CProof}{拓展证明：$y=\frac{k}{x}$ 两支间最短距离（四种方法）}}
统一设 $k>0$，$P(a,\frac{k}{a})$、$Q(-b,-\frac{k}{b})$，$a,b>0$。记 $d=PQ$。

\begin{minipage}{\linewidth}
\par\medskip\textbf{方法一（门槛最低：两次基本不等式）}
\begin{center}
\begin{minipage}{0.98\linewidth}
\centering
\textbf{先分清水平差、竖直差与真实距离}\par\medskip
\begin{tikzpicture}[font=\small,>=Stealth]

\begin{scope}[x=1.65cm,y=1.65cm]
\draw[->,gray] (-1.7,0)--(2.7,0) node[right] {$x$};
\draw[->,gray] (0,-1.6)--(0,1.7) node[above] {$y$};
\draw[CProof,very thick] (-1,-1)--(2,0.5);
\draw[CExample,thick] (-1,-1)--(2,-1)--(2,0.5);
\draw (1.85,-1)--(1.85,-0.85)--(2,-0.85);
\fill (-1,-1) circle (1.7pt) (2,0.5) circle (1.7pt);
\node[above right] at (2,0.5) {$P(a,k/a)$};
\node[below left] at (-1,-1) {$Q(-b,-k/b)$};
\node[below] at (0.5,-1) {$u=a+b$};
\node[right] at (2,-0.45) {$v=k/a+k/b$};
\node[above,rotate=26.565] at (0.6,0.25) {$d=PQ$};
\end{scope}
\end{tikzpicture}\par\smallskip
图取 $k=1,a=2,b=1$。两条直角边为坐标差，斜边才是距离；$d^2\ge2uv$ 的取等要求是 $u=v$。
\end{minipage}
\end{center}

\end{minipage}

\begin{enumerate}[leftmargin=*,itemsep=0.35em]
\item 由两点间距离公式，
\[
d^2=(a+b)^2+\left(\frac{k}{a}+\frac{k}{b}\right)^2.
\]
\item 两个括号均为正数，因此由 $u^2+v^2\ge2uv$，
\[
\begin{aligned}
d^2&\ge 2(a+b)\left(\frac{k}{a}+\frac{k}{b}\right)\\
&=2k\left(2+\frac{a}{b}+\frac{b}{a}\right).
\end{aligned}
\]
\item 因 $\frac{a}{b}>0$、$\frac{b}{a}>0$，且二者乘积为 $1$，故
\[
\frac{a}{b}+\frac{b}{a}\ge2,
\qquad d^2\ge8k.
\]
\item 所以 $d\ge2\sqrt{2k}$。第二次取等需 $a=b$；第一次取等需 $a+b=\frac{k}{a}+\frac{k}{b}$。合起来得到 $a=b=\sqrt{k}$，此时两点确在曲线上，最小值可取到。
\end{enumerate}

\begin{minipage}{\linewidth}
\par\medskip\textbf{方法二（速度最快：对勾函数 $x+\frac{k}{x}$）}
\begin{center}
\begin{minipage}{0.98\linewidth}
\centering
\textbf{斜线段投到 $y=x$ 方向，长度只会缩短}\par\medskip
\begin{tikzpicture}[font=\small,>=Stealth]

\begin{scope}[x=1.6cm,y=1.6cm]
\draw[->,gray] (-0.3,0)--(3.2,0) node[right] {水平差};
\draw[->,gray] (0,-0.3)--(0,2.8) node[above] {竖直差};
\draw[dashed] (0,0)--(2.6,2.6) node[above right] {$y=x$ 方向};
\draw[CProof,very thick] (0,0)--(3,1) node[right] {$(u,v)$};
\draw[CExample,very thick] (0,0)--(2,2);
\draw[dotted,thick] (3,1)--(2,2);
\draw (2.12,1.88)--(2,1.76)--(1.88,1.88);
\node[below] at (1.6,0.4) {$d$};
\node[above left,rotate=45] at (1.45,1.45) {$(u+v)/\sqrt2$};
\node[below left] at (0,0) {$O$};
\end{scope}
\end{tikzpicture}\par\smallskip
平移线段使起点到原点，图示 $u=3,v=1$。投影终点为 $((u+v)/2,(u+v)/2)$；只有 $u=v$ 时，投影长等于 $d$。
\end{minipage}
\end{center}

\end{minipage}

\begin{enumerate}[leftmargin=*,itemsep=0.35em]
\item 记 $u=a+b$、$v=\frac{k}{a}+\frac{k}{b}$，显然 $u,v>0$。
\item 用 $u^2+v^2\ge\frac12(u+v)^2$ 与 $t+\frac{k}{t}\ge2\sqrt{k}$（$t>0$），可得
\[
\begin{aligned}
d&=\sqrt{u^2+v^2}\ge\frac{u+v}{\sqrt2}\\
&=\frac{(a+k/a)+(b+k/b)}{\sqrt2}\\
&\ge\frac{4\sqrt{k}}{\sqrt2}=2\sqrt{2k}.
\end{aligned}
\]
\item 当 $a=b=\sqrt{k}$ 时，两次取等，故即为最小距离。
\end{enumerate}
\par\textbf{意义：} 先优化每一个点的坐标和，再用勾股关系将“坐标和”还原成真实距离；这是从本节对勾函数结论迁移到几何题的最快路线。

\begin{minipage}{\linewidth}
\par\medskip\textbf{方法三（步骤最少：两点坐标先合并）}
\begin{center}
\begin{minipage}{0.98\linewidth}
\centering
\textbf{先固定 $p=ab$，让两点的正参数相等}\par\medskip
\begin{tikzpicture}[font=\small,>=Stealth]

\begin{scope}[x=1.6cm,y=1.3cm]
\draw[->] (0,0)--(4,0) node[right] {$a$};
\draw[->] (0,0)--(0,3.8) node[above] {$b$};
\draw[CProof,very thick,domain=0.3:3.7,samples=100] plot (\x,{1/\x});
\draw[CExample,thick] (0,2)--(2,0);
\draw[dashed] (0,0)--(2.7,2.7) node[right] {$a=b$};
\fill[CProof] (1,1) circle (2.5pt);
\node[above right,fill=white] at (1.2,1.5) {$(\sqrt p,\sqrt p)$};
\node[right] at (2.9,0.6) {$ab=p$};
\node[below] at (2,0) {$2\sqrt p$};
\node[left] at (0,2) {$2\sqrt p$};
\end{scope}
\end{tikzpicture}\par\smallskip
图取 $p=1$。双曲线上 $a+b$ 在切线 $a+b=2\sqrt p$ 处最小，故 $(a+b)^2\ge4p$；再对 $p+k^2/p$ 取最小值，要求 $p=k$。
\end{minipage}
\end{center}

\end{minipage}

\begin{enumerate}[leftmargin=*,itemsep=0.35em]
\item 提取共同因子：
\[
\frac{k}{a}+\frac{k}{b}=\frac{k(a+b)}{ab},
\quad d=(a+b)\sqrt{1+\frac{k^2}{a^2b^2}}.
\]
\item 因 $(a+b)^2\ge4ab$，得
\[
\begin{aligned}
d^2&=(a+b)^2\left(1+\frac{k^2}{a^2b^2}\right)\\
&\ge4\left(ab+\frac{k^2}{ab}\right)\ge8k.
\end{aligned}
\]
最后一个不等式是对正数 $ab$、$\frac{k^2}{ab}$ 用基本不等式。
\item 等号要求 $a=b$ 且 $ab=k$，因此 $a=b=\sqrt{k}$。
\end{enumerate}

\begin{minipage}{\linewidth}
\par\medskip\textbf{方法四（常规几何：两条平行直线夹住曲线）}
\begin{center}
\begin{minipage}{0.98\linewidth}
\centering
\textbf{两支被两条平行线隔开：最近连线垂直穿过}\par\medskip
\begin{tikzpicture}[font=\small,>=Stealth]

\begin{scope}[x=1.05cm,y=1.05cm]
\fill[CProof!6] (-3,1)--(-1,3)--(3,-1)--(1,-3)--cycle;
\draw[->,gray] (-3.3,0)--(3.4,0) node[right] {$x$};
\draw[->,gray] (0,-3.3)--(0,3.4) node[above] {$y$};
\draw[CStatement,very thick,domain=0.32:3.1,samples=100] plot (\x,{1/\x});
\draw[CStatement,very thick,domain=-3.1:-0.32,samples=100] plot (\x,{1/\x});
\draw[CProof,thick] (-1,3)--(3,-1);
\draw[CProof,thick] (-3,1)--(1,-3);
\draw[CExample,very thick] (-1,-1)--(1,1);
\fill[CExample] (1,1) circle (2.3pt) (-1,-1) circle (2.3pt);
\node[right,fill=white] at (1.1,1.1) {$(1,1)$};
\node[left,fill=white] at (-1.1,-1.1) {$(-1,-1)$};
\node[text=CProof] at (2.1,2.1) {$x+y=2$};
\node[text=CProof] at (-2.1,-2.1) {$x+y=-2$};
\draw (0.86,0.86)--(0.72,1)--(0.86,1.14);
\end{scope}
\end{tikzpicture}\par\smallskip
图取 $k=1$，强调线段长 $2\sqrt2$。一般情形：两条线为 $x+y=\pm2\sqrt k$，间距 $2\sqrt{2k}$，最近点为 $(\pm\sqrt k,\pm\sqrt k)$。
\end{minipage}
\end{center}

\end{minipage}

\begin{enumerate}[leftmargin=*,itemsep=0.35em]
\item 对第一象限的任意点 $P(x,y)$，有 $xy=k$，于是
\[
x+y\ge2\sqrt{xy}=2\sqrt{k}.
\]
第一象限分支全部位于直线 $\ell_1:x+y=2\sqrt{k}$ 的同一侧，且在 $(\sqrt{k},\sqrt{k})$ 与这条直线相接。
\item 对第三象限的任意点 $Q(x,y)$，令 $u=-x>0,v=-y>0$，则 $uv=k$，
\[
-(x+y)=u+v\ge2\sqrt{k},\quad x+y\le-2\sqrt{k}.
\]
第三象限分支位于 $\ell_2:x+y=-2\sqrt{k}$ 的另一侧，且在 $(-\sqrt{k},-\sqrt{k})$ 与该直线相接。
\item 因两条直线平行，间距为
\[
\frac{|2\sqrt{k}-(-2\sqrt{k})|}{\sqrt{1^2+1^2}}
=2\sqrt{2k}.
\]
处在两侧的 $P,Q$ 的距离不小于该间距；取上述两个相接点时，连线垂直于两条直线，刚好取得这个距离。
\end{enumerate}

\textbf{结论：} 四法均得 $\boxed{d_{\min}=2\sqrt{2k}}$，当且仅当 $P=(\sqrt{k},\sqrt{k})$、$Q=(-\sqrt{k},-\sqrt{k})$。特别地 $k=1$ 时 $d_{\min}=2\sqrt2$。

\end{proofbox}
